% Generated draft: requires mathematical and source review.
\documentclass[11pt]{article}
\usepackage[margin=25mm]{geometry}
\usepackage{fontspec}
\setmainfont{Latin Modern Roman}
\usepackage{amsmath,amssymb,mathtools,unicode-math}
\setmathfont{Latin Modern Math}
\usepackage{enumitem,needspace,xurl,hyperref}
\setlength{\parindent}{0pt}
\setlength{\parskip}{5pt}
\setlength{\emergencystretch}{3em}
\newcommand{\N}{\mathbb{N}}
\newcommand{\Z}{\mathbb{Z}}
\newcommand{\Q}{\mathbb{Q}}
\newcommand{\R}{\mathbb{R}}
\newcommand{\C}{\mathbb{C}}
\newcommand{\F}{\mathbb{F}}
\newcommand{\A}{\mathbb{A}}
\newcommand{\PP}{\mathbb{P}}
\newcommand{\E}{\mathbb{E}}
\newcommand{\eps}{\varepsilon}
\newcommand{\dd}{\,\mathrm{d}}
\begin{document}
\section*{UnsolvedMath ID: 30005280}
\textbf{Source problem code:} catalogue entry
\subsection{Definitions and mathematical statement}
The supplied registry labels this problem \emph{Onsager Thresholds for Hydrostatic Euler Energy Conservation}. Its stated target is to determine critical anisotropic regularity thresholds for conservation of horizontal kinetic energy by weak solutions of the hydrostatic Euler equations, and to identify lower regularities at which anomalous dissipation can occur.

For purposes of explaining the energy terminology only, if a spatial domain \(\Omega\), a time interval \(I\), and a horizontal velocity field \(u_h\) have been specified, with \(u_h(t,\cdot)\in L^2(\Omega)\) for almost every \(t\in I\), define the horizontal kinetic energy by
\[
E_h(t)=\frac12\int_{\Omega}|u_h(t,x)|^2\,\dd x.
\]
Conservation means that there exists a constant \(E_0\) such that \(E_h(t)=E_0\) for almost every \(t\in I\). Here anisotropic regularity refers to regularity conditions that distinguish horizontal and vertical spatial directions; the record specifies neither the function spaces nor their exponents.

\textbf{Requested mathematical target, with unresolved qualifications.} Determine conditions on anisotropic regularity under which every weak solution in the intended class conserves horizontal kinetic energy, and determine the lower-regularity regime in which weak solutions exhibiting anomalous dissipation exist. This is a restatement of the supplied target, not a fully specified conjecture: the equations, spatial dimension and domain, boundary conditions, forcing assumptions, time interval, weak-solution class, integrability requirements, regularity scale, and meaning of a critical threshold are absent. The supplied text also does not specify whether dissipation is required at every regularity below a threshold, or only at some such regularities, nor what is asserted at the threshold itself. To identify energy loss as anomalous dissipation requires specifying the energy balance, including any forcing and boundary fluxes.

The original extract is a table-of-contents fragment beginning with an unattached word ``conjecture'' and continuing with titles of other workshop contributions. It does not supply the hypotheses needed to complete the requested statement. The registry's claim that the corrected statement was verified against the report conflicts with its supplied literature assessment that no definite conjecture or proposition is identifiable. These qualifications remain unresolved; no precise threshold claim is asserted here.
\subsection{Short English statement}
The supplied target asks how much horizontal and vertical regularity guarantees horizontal kinetic-energy conservation for weak hydrostatic Euler solutions, and at what lower regularities dissipative solutions can exist. The supplied extract does not specify the equations, solution class, boundary conditions, or regularity spaces, so a precise threshold conjecture cannot be recovered. The recorded open status and verification claims are supplied curation and require human review.
\subsection{Sources}
\begin{itemize}
\item Local registry: \texttt{data/unsolvedmath/problems.json}, permanent UnsolvedMath ID \(30005280\), exact source problem code \texttt{catalogue entry}; supplied complete source record. The record gives proposed year \(2022\) and supplies no named proposer.
\item Supplied report citation: \emph{Mathematical Advances in Geophysical Fluid Dynamics} (2023), Oberwolfach Reports, DOI 10.4171/owr/2022/51, pp.\ 2961--3003. \url{https://doi.org/10.4171/owr/2022/51}.
\item Supplied literature source: EMS Press, OWR 2022/51. \url{https://doi.org/10.4171/owr/2022/51}. The exact supplied literature-check date is \texttt{2026-08-22}. This assessment is supplied dated curation, not independently verified: it classifies the record as open because the extracted table-of-contents fragment does not identify a definite conjecture or proposition. The supplied background requests independent verification of time-sensitive status and cited claims.
\item The registry marks the statement \texttt{corrected\_verified} and says it was checked against the official report. That supplied verification claim is not independently established here and is in tension with the supplied literature assessment. The report-level citation does not identify a supporting passage or author for this particular hydrostatic Euler target; names and joint-work credits in the fragment concern other listed contributions and cannot establish its attribution.
\end{itemize}
\end{document}
